\scanpage{4}
\mspage{2}
\uline{Ex}\qquad $C:\mathrm{Modf}(k)$\qquad
$F$ est le foncteur $M\otimes_k k'$.

Alors
\[
U=k'\otimes_k k'
\]
(structure de $k'$-alg. provenant du facteur gauche),
section canonique $\varepsilon:k'\otimes k'\longrightarrow k'$,
donc
\[
U^{\vee(k')}=A\simeq k'\otimes_k {k'}^{\vee(k)}
\simeq\operatorname{End}_k(k'),
\]
et on y prend la structure d'anneau habituelle.

\margin{\uline{NB} \unsure{alors}
\inserted{(si $k'/k$ est \unsure{ét.|étale}, donc $U\simeq A^{\vee(k)}$)}
$U\xrightarrow{\Delta_U}U\otimes_k U$ transposé de
$A\otimes_k A\longrightarrow A$ \unsure{est une}
\deleted{\unsure{structure}} \unsure{coalg.} \uline{\unsure{co-}}associative
($\Delta_U(1)\ne1$ en général, et $\Delta_U$ pas multiplicatif\ldots)}

\par\noindent\rule{.85\linewidth}{.4pt}\par
On cherche conditions \inserted{intrinsèques} sur \inserted{la $\otimes$-cat.} $C$ qui
assurent que pour \deleted{\unsure{$A\in C$}}
$k'/k$ et $F$ comme dessus, le foncteur des
$A_{\unsure{k'|k}}\inserted{=A\otimes_k k'}$-modules de \unsure{t.f.|d.f.}
\[
M'\longrightarrow M'\otimes_{k'\otimes_k k'}k'
\]
soit \uline{exact} (OK si $k'/k$ est étale) et \uuline{fidèle}.

Suffit il que $C$ soit \uline{rigide}, et que
\[
k\xrightarrow{\sim}\operatorname{End}_C(1_C)\quad???
\]
\par\hspace{.3\linewidth}\rule{.25\linewidth}{.4pt}\par
\uline{Exprimer la rigidité} de $C$ !
\margin{3}
